% from ega2-1-fr.tex
\oldpage[II]{63}

et que $X$ est de type fini sur $Y$~; il en résulte alors que $\mathcal{S}'$
est engendrée par $\mathcal{S}'_1$ (comme on le voit en se ramenant au cas où
$Y$ et $Y'$ sont affines) et que $X'$ est de type fini sur $Y'$
(\hyperref[I.6.3.4-fr]{I, 6.3.4}). Soit $\mathcal{F}$ un
$\mathcal{O}_X$-Module et posons $\mathcal{F}'=\Psi^*(\mathcal{F})$~; il
résulte alors de (\hyperref[II.3.5.4-fr]{3.5.4}) et de
(\hyperref[0.4.3.3-fr]{0, 4.3.3}) que l'on a
$\mathcal{F}'(n)=\Psi^*(\mathcal{F}(n))$ pour tout $n\in\mathbf{Z}$. On
définit en outre un $\psi$-homomorphisme canonique
$\theta_n:q_*(\mathcal{F}(n))\to q'_*(\mathcal{F}'(n))$ de la façon
suivante~: vu la commutativité du diagramme
\[
  \xymatrix{
    X\ar[d]_{q}
    & X'\ar[l]_{\Psi}\ar[d]^{q'}\\
    Y
    & Y'\ar[l]^{\psi}
  }
\]
il s'agit de définir un homomorphisme
$q_*(\mathcal{F}(n))\to
\psi_*(q'_*(\Psi^*(\mathcal{F}(n))))
=q_*(\Psi_*(\Psi^*(\mathcal{F}(n))))$, et il suffit de prendre
l'homomorphisme $\theta_n=q_*(\rho_n)$, où $\rho_n$ est l'homomorphisme
canonique
$\mathcal{F}(n)\to\Psi_*(\Psi^*(\mathcal{F}(n)))$
(\hyperref[0.4.4.3-fr]{0, 4.4.3}). Il est immédiat que pour tout ouvert affine
$U$ de $Y$ et tout ouvert affine $U'$ de $Y'$ tel que
$U'\subset\psi^{-1}(U)$, l'homomorphisme $\theta_n$ donne pour les sections
l'homomorphisme canonique (\hyperref[0.3.7.2-fr]{0, 3.7.2})
$\Gamma(q^{-1}(U),\mathcal{F}(n))\to
\Gamma(q'^{-1}(U'),\mathcal{F}'(n))$.

La commutativité de (\hyperref[II.2.8.13.2-fr]{2.8.13.2}) montre alors que si
$\mathcal{F}$ est quasi-cohérent, le diagramme
\[
  \xymatrix{
    \mathcal{F}\ar[r]^{\rho}
    & \mathcal{F}'\\
    (\Gamma_*(\mathcal{F}))^{\sim}
      \ar[u]^{\beta_{\mathcal{F}}}\ar[r]_{\widetilde{\theta}}
    & (\Gamma_*(\mathcal{F}'))^{\sim}
      \ar[u]_{\beta_{\mathcal{F}'}}
  }
\]
est commutatif (la flèche horizontale du haut étant le $\Psi$-morphisme
canonique $\mathcal{F}\to\Psi^*(\mathcal{F})$).

De même, la commutativité de
(\hyperref[II.2.8.13.1-fr]{2.8.13.1}) montre que le diagramme
\[
  \xymatrix{
    \Gamma_*(\widetilde{\mathcal{M}})\ar[r]^{\theta}
    & \Gamma_*(\widetilde{\mathcal{M}'})\\
    \mathcal{M}\ar[r]_{\rho}\ar[u]^{\alpha_{\mathcal{M}}}
    & \mathcal{M}'\ar[u]_{\alpha_{\mathcal{M}'}}
  }
\]
est commutatif (la flèche horizontale du bas étant le $\psi$-morphisme
canonique $\mathcal{M}\to\psi^*(\mathcal{M})$).
\end{env}

\begin{env}[3.5.6]
\label{II.3.5.6-fr}
Considérons maintenant deux préschémas $Y$, $Y'$, un morphisme
$g:Y'\to Y$, une $\mathcal{O}_Y$-Algèbre (resp.
$\mathcal{O}_{Y'}$-Algèbre) graduée quasi-cohérente $\mathcal{S}$ (resp.
$\mathcal{S}'$) et un $g$-morphisme d'Algèbres graduées
$u:\mathcal{S}\to\mathcal{S}'$, c'est-à-dire un
$\mathcal{O}_Y$-homomorphisme d'Algèbres graduées
$\mathcal{S}\to g_*(\mathcal{S}')$~; on sait d'ailleurs qu'il revient au même
de se donner un $\mathcal{O}_{Y'}$-homomorphisme d'Algèbres graduées
$u^\sharp:g^*(\mathcal{S})\to\mathcal{S}'$. On déduit alors canoniquement de
$u^\sharp$ un $Y'$-morphisme
$W=\operatorname{Proj}(u^\sharp):G(u^\sharp)\to
\operatorname{Proj}(g^*(\mathcal{S}))$ où $G(u^\sharp)$ est un ouvert de
$X'=\operatorname{Proj}(\mathcal{S}')$
(\hyperref[II.3.5.1-fr]{3.5.1}). D'autre part,
$X''=\operatorname{Proj}(g^*(\mathcal{S}))$ s'identifie

